\section{Research Experience}

    \begin{twocolentry}{Washington, DC \\ Jun 2026 to Aug 2026}
        \textbf{Naval Research Laboratory (NRL)}\\
        \textit{NREIP Intern, Laboratory for Autonomous Systems Research (5517)}
    \end{twocolentry}

    \vspace{0.10 cm}
    \begin{onecolentry}
        \begin{highlights}
        	\item Analyzed ERA5 and radiosonde atmospheric data to support research on optimal station-keeping locations for solar high-altitude balloons (sHABs)
            \item Redesigned the forecast processing pipeline by unifying ERA5 and GFS workflows, adding multi-year data support, and improving pipeline reliability through resumable multiprocessing in Python
            \item Increased the accuracy of wind analysis by optimizing interpolation methods and updating Hovmöller wind-direction processing, while developing tools to validate and migrate historical datasets
            %\item Programmed and maintained solar high-altitude balloon systems for upcoming research flights and atmospheric data collection.
        \end{highlights}
    \end{onecolentry}  
    
    \vspace{0.10 cm}
    
    \begin{twocolentry}{College Park, MD \\ Aug 2025 to Present}
         \href{https://ottelab.com}{\textbf{University of Maryland Motion and Teaming Laboratory}}\\
        \textit{Graduate Research Assistant}
    \end{twocolentry}

    \vspace{0.10 cm}
    \begin{onecolentry}
        \begin{highlights}
        	\item Analyzing Underwater SLAM applications for multi-vehicle navigation systems and current methodologies and identifying research gaps for autonomous underwater vehicle coordination
        \end{highlights}
    \end{onecolentry}  
    
    \vspace{0.20 cm}
    
    \begin{twocolentry}{Bridgeport, CT \\ Dec 2023 to Aug 2025}
        \href{https://nustem.bridgeport.edu/}{\textbf{Exploring Naval Underwater STEM (Exploring NUSTEM)}\\ 
        \textit{Undergraduate Researcher}}
    \end{twocolentry}

    \vspace{0.10 cm}
    \begin{onecolentry}
        \begin{highlights}
            \item Constructed and modified three BlueROV2 underwater vehicle platforms with integrated sonar systems, supporting successful data collection at depths of approximately 300 m in Long Island Sound
            \item Diagnosed electrical and material degradation in the Explorer submarine using oscilloscope-based testing and component inspection, then implemented corrosion remediation and protective treatments to support restoration and system reliability
            %\item Led technical design reviews with industry professionals during restoration of the Explorer submarine, evaluating subsystem integration, operational readiness, and safety requirements for future deployment
            %\item Funded under Office of Naval Research (ONR) FOA N00014-21-S-F004 Science, Technology, Engineering \& Mathematics (STEM), Education and Workforce Program, “Exploring Naval Underwater STEM (Exploring NUSTEM)”
        \end{highlights}
    \end{onecolentry}

    \vspace{0.10 cm}

    \begin{twocolentry}{Bridgeport, CT \\ Sept 2022 to Aug 2025}
        \href{http://eclipse.montana.edu/}{\textbf{Nationwide Eclipse Ballooning Project (NEBP)}}\\ 
        \textit{Student Pod Lead/ Systems Engineering Lead}
    \end{twocolentry}

    \vspace{0.10 cm}
    \begin{onecolentry}
        \begin{highlights}
            \item Led 10 engineers in designing and qualifying high-altitude balloon payload systems that collected eclipse data in 100+ mph winds at altitudes above 100,000 ft
            \item Coordinated hardware and software integration across student and faculty teams in the East Central Region, resolving subsystem-level issues and ensuring completion of project milestones for multi-site eclipse missions
            \item Coordinated the standardization of datasets from 53 participating teams by removing non-flight data, labeling valid flight intervals, and normalizing timestamps to UTC, supporting cross-team correlation and publication of the 2023–2024 eclipse campaign dataset on Dryad \href{https://datadryad.org/dataset/doi:10.5061/dryad.rv15dv4jh}{(Des Jardins, et al. 2025)}
            %\item Funded under National Aeronautics And Space Administration, Subaward Title: Montana Space Grant Consortium Opportunities in STEM 2020-2024, Subaward ID: G355-22-W9649
        \end{highlights}
    \end{onecolentry}
\begin{comment}
	    \vspace{0.2 cm}

    \begin{twocolentry}{
        Apr 2022 to Sept 2023
    }
        \href{https://www.glee2023.org}{\textbf{Team 2149 Lead}, Great Lunar Expedition for Everyone (GLEE) -- Bridgeport, CT}
    \end{twocolentry}

    \vspace{0.15 cm}
    \begin{onecolentry}
        \begin{highlights}
          \item Led an interdisciplinary team of 10 to design, test, and prepare for the launch of “LunaSat” payloads, which were engineered to conduct localized scientific missions on the lunar surface
          \item Facilitated team meetings and cross-functional collaboration, improving communication and ensuring efficient problem-solving throughout the project
        \end{highlights}
    \end{onecolentry}
\end{comment}